\documentclass[11pt,a4paper]{article}
\usepackage[T1]{fontenc}
\usepackage[utf8]{inputenc}
\usepackage{amsmath,amssymb,amsthm}
\usepackage[margin=1in]{geometry}
\usepackage[colorlinks=true,linkcolor=blue,urlcolor=blue]{hyperref}
\theoremstyle{plain}
\newtheorem{theorem}{Theorem}
\newtheorem{lemma}[theorem]{Lemma}
\newtheorem{proposition}[theorem]{Proposition}
\newtheorem{corollary}[theorem]{Corollary}
\theoremstyle{definition}
\newtheorem{remark}[theorem]{Remark}
% A quoted conjecture can be a single hundred-term formula whose terms are thirty-digit
% numbers. TeX cannot hyphenate those, so without a little stretch every such quote runs off
% the page; the linked-a-file papers are all of that shape.
\setlength{\emergencystretch}{3em}

\title{The empirical closed form for OEIS A186088, proved}
\author{Adrian Perez Fontelles\\ \small Independent researcher}
\date{1 September 2026}
\begin{document}
\maketitle

\begin{abstract}
OEIS A186088 carries an empirical closed form for $a(n)$ --- not a recurrence relating
consecutive terms but an explicit formula. It is true, and it is decidable rather than
empirical. The entry counts arrays subject to a condition on a bounded window of consecutive lines, so the lines of an array are the
vertices of a finite digraph, an array is a walk in it, and $a(n)$ is a walk count on
$S=291$ vertices. Any function of the form $\sum_j p_j(n)\lambda_j^{\,n}$ satisfies the
monic linear recurrence whose characteristic polynomial is
$\prod_j(x-\lambda_j)^{\deg p_j+1}$; here that polynomial is $q(x)=(x-1)^{35}$, of degree
$35$. Two things are then checked exactly: that the walk count satisfies the same
recurrence from some index on, which the Cayley--Hamilton theorem reduces to a finite
computation, and that the two sequences agree at $35$ consecutive indices past that
point. A monic recurrence determines every later term from its predecessors, so agreement
there is agreement for good.
\end{abstract}

\noindent\small 2020 Mathematics Subject Classification. 05A15, 11B37, 15A18.\normalsize

\section{The conjecture}

OEIS A186088 is ``Number of (n+2)X3 0..4 arrays with each 3X3 subblock having rows and columns in lexicographically nondecreasing order.''. It has offset $1$ and begins
\[
102251, 1252889, 11258613, 83378583, 531218757, 2985984444, 15084070635, 69482992431,\ \dots
\]
The entry states, as a conjecture contributed by its author and never marked settled:
\begin{quote}\raggedright\small
Empirical: a(n) = (1/295232799039604140847618609643520000000)*n\^{}34 + (1/36561337342365837875866081689600000)*n\^{}33 + (1151/78939251080108059050165403648000000)*n\^{}32 + (57793/15180625207713088278877962240000000)*n\^{}31 + (8146051/12732137270985170814542807040000000)*n\^{}30 + (15252793/195879034938233397146812416000000)*n\^{}29 + (917629/125082397789421070974976000000)*n\^{}28 + (40844936773/73173202706811326520360960000000)*n\^{}27 + (10326664514897/292692810827245306081443840000000)*n\^{}26 + (1699793021071/900593264083831711019827200000)*n\^{}25 + (21419403120091/247663147623053720530452480000)*n\^{}24 + (3257433593147959/952550567780975848194048000000)*n\^{}23 + (149148353528547551/1272577438379327417745408000000)*n\^{}22 + (888681424414566953/254515487675865483549081600000)*n\^{}21 + (11531836957656612161/127257743837932741774540800000)*n\^{}20 + (5909972330146165909/2879134475971328999424000000)*n\^{}19 + (185196564663366054480883/4554487674199698126667776000000)*n\^{}18 + (223744365124534317893/317054484803320440422400000)*n\^{}17 + (631124705896508562289697/58734343309815111588249600000)*n\^{}16 + (7215220158705833717523107/50233319936026082279424000000)*n\^{}15 + (44200951728497459126549888033/26246909666573627990999040000000)*n\^{}14 + (1394814879695359558279349521/80759722050995778433843200000)*n\^{}13 + (32548509299450082036607825951/211076546269648057270272000000)*n\^{}12 + (48337137666188129039919297713/40591643513393857167360000000)*n\^{}11 + (18053107623124660173941442743743/2286662584587853953761280000000)*n\^{}10 + (2266224377869928249871300899693/50814724101952310083584000000)*n\^{}9 + (902210285827865272553906095027/4234560341829359173632000000)*n\^{}8 + (57964514565563064206413856761/67861543939573063680000000)*n\^{}7 + (2023140467628393995969236886677/710661168478306805760000000)*n\^{}6 + (25417310330253175708558630751/3279974623746031411200000)*n\^{}5 + (61858293620765291251250539109/3728237822324655704064000)*n\^{}4 + (12364083190538430205351213/471396394434917952000)*n\^{}3 + (8255316109684330210707767/294877831150846944000)*n\^{}2 + (1418176238189177/80224196052)*n + 2040
\end{quote}
As of the ``Last modified'' line on the live entry (Oct 03 2025 03:24:28, revision 9) this is still
recorded as empirical, and nothing on the entry records it as proved. Write
\begin{multline*}
f(n)\;=\;\frac{n^{34}}{295232799039604140847618609643520000000} + \frac{n^{33}}{36561337342365837875866081689600000} \\
+ \frac{1151 n^{32}}{78939251080108059050165403648000000} + \frac{57793 n^{31}}{15180625207713088278877962240000000} \\
+ \frac{8146051 n^{30}}{12732137270985170814542807040000000} + \frac{15252793 n^{29}}{195879034938233397146812416000000} \\
+ \frac{917629 n^{28}}{125082397789421070974976000000} + \frac{40844936773 n^{27}}{73173202706811326520360960000000} \\
+ \frac{10326664514897 n^{26}}{292692810827245306081443840000000} + \frac{1699793021071 n^{25}}{900593264083831711019827200000} \\
+ \frac{21419403120091 n^{24}}{247663147623053720530452480000} + \frac{3257433593147959 n^{23}}{952550567780975848194048000000} \\
+ \frac{149148353528547551 n^{22}}{1272577438379327417745408000000} + \frac{888681424414566953 n^{21}}{254515487675865483549081600000} \\
+ \frac{11531836957656612161 n^{20}}{127257743837932741774540800000} + \frac{5909972330146165909 n^{19}}{2879134475971328999424000000} \\
+ \frac{185196564663366054480883 n^{18}}{4554487674199698126667776000000} + \frac{223744365124534317893 n^{17}}{317054484803320440422400000} \\
+ \frac{631124705896508562289697 n^{16}}{58734343309815111588249600000} + \frac{7215220158705833717523107 n^{15}}{50233319936026082279424000000} \\
+ \frac{44200951728497459126549888033 n^{14}}{26246909666573627990999040000000} + \frac{1394814879695359558279349521 n^{13}}{80759722050995778433843200000} \\
+ \frac{32548509299450082036607825951 n^{12}}{211076546269648057270272000000} + \frac{48337137666188129039919297713 n^{11}}{40591643513393857167360000000} \\
+ \frac{18053107623124660173941442743743 n^{10}}{2286662584587853953761280000000} + \frac{2266224377869928249871300899693 n^{9}}{50814724101952310083584000000} \\
+ \frac{902210285827865272553906095027 n^{8}}{4234560341829359173632000000} + \frac{57964514565563064206413856761 n^{7}}{67861543939573063680000000} \\
+ \frac{2023140467628393995969236886677 n^{6}}{710661168478306805760000000} + \frac{25417310330253175708558630751 n^{5}}{3279974623746031411200000} \\
+ \frac{61858293620765291251250539109 n^{4}}{3728237822324655704064000} + \frac{12364083190538430205351213 n^{3}}{471396394434917952000} \\
+ \frac{8255316109684330210707767 n^{2}}{294877831150846944000} + \frac{1418176238189177 n}{80224196052} \\
+ 2040.
\end{multline*}

\section{The count is a walk count}

The entry's defining condition is local: it is a condition on a bounded window of consecutive lines. Take as vertices the
admissible configurations of such a window and put an edge from $u$ to $v$ when $v$ may
follow $u$, which the condition decides by inspection. An admissible array is then exactly a
walk, so with $M$ the adjacency matrix of that digraph on $S=291$ vertices, and $u,w$ the
vectors recording which windows may start and end an array,
\[
a(n)\;=\;u^{\!\top}M^{\,g(n)}w
\]
for the linear function $g$ that the entry's shape supplies. In particular $a$ is
$C$-finite: it satisfies some monic integer linear recurrence, though not necessarily a short
one.

\section{A closed form is a recurrence}

\begin{lemma}\label{lem:ann}
Let $f(m)=\sum_{j}p_j(m)\lambda_j^{\,m}$ with the $\lambda_j$ distinct and the $p_j$
polynomials, and put $q(x)=\prod_j(x-\lambda_j)^{\deg p_j+1}$. Then $q$ applied to the shift
operator annihilates $f$: writing $q(x)=x^{r}-\sum_{i=1}^{r}c_ix^{r-i}$,
\[
f(m)\;=\;\sum_{i=1}^{r}c_i\,f(m-i)\qquad\text{for every }m.
\]
\end{lemma}

\begin{proof}
The shift operator $E$ acts on $m\mapsto\lambda^{m}p(m)$ as
$(E-\lambda)\bigl(\lambda^{m}p(m)\bigr)=\lambda^{m+1}\bigl(p(m+1)-p(m)\bigr)$, and
$p(m+1)-p(m)$ has degree one less than $p$. So $(E-\lambda)^{\deg p+1}$ kills
$\lambda^{m}p(m)$, and the factors of $q$ commute.
\end{proof}

Here $q(x)=(x-1)^{35}$, of degree $35$; the closed form is a polynomial in $n$, so this is $(x-1)^{35}$, and the recurrence it encodes is
\begin{equation}\label{eq:rec}
a(n)\;=\;35\,a(n-1) - 595\,a(n-2) + 6545\,a(n-3) - 52360\,a(n-4) + 324632\,a(n-5) - 1623160\,a(n-6) + 6724520\,a(n-7) - 23535820\,a(n-8) + 70607460\,a(n-9) - 183579396\,a(n-10) + 417225900\,a(n-11) - 834451800\,a(n-12) + 1476337800\,a(n-13) - 2319959400\,a(n-14) + 3247943160\,a(n-15) - 4059928950\,a(n-16) + 4537567650\,a(n-17) - 4537567650\,a(n-18) + 4059928950\,a(n-19) - 3247943160\,a(n-20) + 2319959400\,a(n-21) - 1476337800\,a(n-22) + 834451800\,a(n-23) - 417225900\,a(n-24) + 183579396\,a(n-25) - 70607460\,a(n-26) + 23535820\,a(n-27) - 6724520\,a(n-28) + 1623160\,a(n-29) - 324632\,a(n-30) + 52360\,a(n-31) - 6545\,a(n-32) + 595\,a(n-33) - 35\,a(n-34) + a(n-35).
\end{equation}

\section{The proof}

\begin{lemma}\label{lem:crit}
With $q$ as above, $a$ satisfies \eqref{eq:rec} for all large $n$ if and only if
$u^{\!\top}M^{\,j}q(M)w=0$ for every $j\ge0$, and it is enough to check $j<S$.
\end{lemma}

\begin{proof}
Substituting the walk expression, the residual $a(n)-\sum_ic_ia(n-i)$ equals
$u^{\!\top}M^{\,g(n)-r}q(M)w$ up to the fixed linear reparametrisation $g$. For the bound,
$M^{S}$ is an integer combination of $I,M,\dots,M^{S-1}$ by Cayley--Hamilton, so the values
$u^{\!\top}M^{j}q(M)w$ satisfy the monic recurrence given by the characteristic polynomial of
$M$; $S$ consecutive zeros therefore force all later ones, and the last nonzero value pins the
threshold exactly.
\end{proof}

\begin{theorem}
$a(n)=f(n)$ for every $n\ge1$.
\end{theorem}

\begin{proof}
The vector $w'=q(M)w$ was formed by $35$ matrix--vector products in exact integer
arithmetic, and $u^{\!\top}M^{j}w'$ evaluated for $j=0,1,\dots$, again exactly. Those integers
vanish from the point corresponding to $n=36$ onwards, and the run was continued until the
working vector was identically zero, which settles every larger $j$ at once. So $a$ satisfies
\eqref{eq:rec} for every $n>35$. By Lemma~\ref{lem:ann} $f$ satisfies \eqref{eq:rec}
identically. Both were then evaluated in exact arithmetic at the $35$ consecutive indices
$n=36,\dots,70$ and agree at each. A monic recurrence of order $35$
determines $a(m)$ from $a(m-1),\dots,a(m-35)$, and for every $m\ge71$ both
$m>35$ and $m-35\ge36$ hold, so induction on $m$ gives $a(m)=f(m)$ for every
$m\ge36$.
Below $n=36$ the recurrence argument gives nothing, since \eqref{eq:rec} is only
established for $a$ from $n=36$ on. The remaining indices $n=1,\dots,35$ were
therefore checked one at a time, directly against the walk count in exact integer arithmetic,
and agree.
\end{proof}

The range is tight in the sense that was checked: the closed form holds from the entry's first term onwards.

\section{Verification}

Four checks, each independent of the algebra above.

First, the walk model reproduces every one of the $16$ terms the entry publishes,
exactly, in integer arithmetic. This is what ties the model to the entry: the model is built
by reading the entry's English, and a misreading gives a different digraph and different
counts.

Second, the closed form was evaluated in exact rational arithmetic --- not floating point ---
at every published term from $n=1$ on, and agrees with the entry's own DATA at each.

Third, the annihilation test of Lemma~\ref{lem:crit} was carried out exactly, and the model
and the closed form were then compared directly at sixty consecutive indices beyond
$n=1$, far past where the proof already settles the matter.

Fourth, the closed form was deliberately perturbed --- by a constant and by a term linear in
$n$ --- and each perturbed form was rejected by the same comparison, so the test is not one
that anything would pass.

\begin{thebibliography}{9}
\bibitem{oeis} The OEIS Foundation, \emph{The On-Line Encyclopedia of Integer
Sequences}, \url{https://oeis.org}, 2026. Sequence A186088.
\bibitem{stanley} R.~P.~Stanley, \emph{Enumerative Combinatorics}, Volume 1, 2nd ed.,
Cambridge University Press, 2011. (Transfer-matrix method, Section 4.7; $C$-finite sequences,
Section 4.1.)
\bibitem{kp} M.~Kauers and P.~Paule, \emph{The Concrete Tetrahedron}, Springer, 2011.
(Closed forms and linear recurrences with constant coefficients.)
\bibitem{horn} R.~A.~Horn and C.~R.~Johnson, \emph{Matrix Analysis}, 2nd ed., Cambridge
University Press, 2013.
\end{thebibliography}

\end{document}
